% GENERATED FILE. Edit canonical lessons, then run npm run generate:books.
% canonical-source-hash: fnv1a64:f1f4a4f1a91485e5
% canonical-lessons: LA-C349-fabulor, LA-C349-aio, LA-C349-mirabilis, LA-C349-incredibilis, LA-C349-commentum

\chapter{Fābulae mīrae}
\label{ch:la-fabulae-mirae}
\hlchaptermodality{349}

\begin{chapteropening}
\textbf{By the end of this chapter:} \emph{I can tell and react to a tall story: to chat, to say, wonderful, unbelievable and a made-up story.}

Complete the last lesson of chapter 349: I can tell and react to a tall story: to chat, to say, wonderful, unbelievable and a made-up story.
\end{chapteropening}

\section[fābulor, fābulārī]{fābulor, fābulārī — to chat, to tell stories}

\label{lesson:LA-C349-fabulor}



\begin{quote}
Before the new one: say the Latin for an omen, then the Latin for to prophesy.
\end{quote}

\subsection*{You'll want to know: fābulor, fābulārī}
\textbf{fābulor, fābulārī} — \textquotedblleft{}to chat, to tell stories\textquotedblright{}.

\begin{quote}
Senēs in tabernā dē rēbus antīquīs fābulābantur.
\end{quote}

\begin{quote}
\emph{The old men were chatting in the inn about the old days.}
\end{quote}

\emph{fābulor, fābulārī, fābulātus sum}, a deponent, from \emph{fābula} (taught). In Spain it became \emph{hablar}, 'to speak'.

\begin{cousinweb}
\textbf{fābulor} is built on \emph{fābula}, \textquotedblleft{}story, talk\textquotedblright{} (Lewis \& Short), from \emph{fārī}, \textquotedblleft{}to speak\textquotedblright{}. Spanish \emph{hablar}, \textquotedblleft{}to speak\textquotedblright{}, comes from it.
\end{cousinweb}

\begin{grammarlens}[title={What you've built}]
One more word for telling, explaining and describing at B1.
\end{grammarlens}

\subsection*{Your turn}
\begin{itemize}
\raggedright
  \item \emph{Say it:} \emph{fābulor, fābulārī}
  \item \emph{Say it:} \emph{fābulor, fābulārī}, once more
  \item \emph{Say it:} a sentence of your own with \emph{fābulor, fābulārī}
  \item \emph{Recall:} say the Latin for an omen, then the Latin for to prophesy, then say \emph{fābulor, fābulārī} again
\end{itemize}

\begin{tcolorbox}[breakable,colback=teal!4,colframe=teal!35!black,title={Before you move on}]
What does fābulor, fābulārī mean? (\textquotedblleft{}To chat, to tell stories\textquotedblright{}.) Say it once more.
\end{tcolorbox}
\section[āiō]{āiō — to say yes; to say}

\label{lesson:LA-C349-aio}



\begin{quote}
Before the new one: say the Latin for to prophesy, then the Latin for to chat.
\end{quote}

\subsection*{You'll want to know: āiō}
\textbf{āiō} — \textquotedblleft{}to say yes; to say\textquotedblright{}.

\begin{quote}
Nauta mōnstrum sē vīdisse ait.
\end{quote}

\begin{quote}
\emph{The sailor says he saw a monster.}
\end{quote}

Defective: mostly \emph{āiō, ais, ait, aiunt}. \emph{Aiunt}: they say. Its opposite is \emph{negō} (taught), 'to say no'.

\begin{cousinweb}
\textbf{āiō} is inherited from Proto-Italic \emph{*agjō} (EtymDB), from a Proto-Indo-European root meaning \textquotedblleft{}to say\textquotedblright{}. Only a few of its forms are in use (Lewis \& Short).
\end{cousinweb}

\begin{grammarlens}[title={What you've built}]
One more word for telling, explaining and describing at B1.
\end{grammarlens}

\subsection*{Your turn}
\begin{itemize}
\raggedright
  \item \emph{Say it:} \emph{āiō}
  \item \emph{Say it:} \emph{āiō}, once more
  \item \emph{Say it:} a sentence of your own with \emph{āiō}
  \item \emph{Recall:} say the Latin for to prophesy, then the Latin for to chat, then say \emph{āiō} again
\end{itemize}

\begin{tcolorbox}[breakable,colback=teal!4,colframe=teal!35!black,title={Before you move on}]
What does āiō mean? (\textquotedblleft{}To say yes; to say\textquotedblright{}.) Say it once more.
\end{tcolorbox}
\section[mīrābilis, mīrābile]{mīrābilis, mīrābile — wonderful, astonishing}

\label{lesson:LA-C349-mirabilis}



\begin{quote}
Before the new one: say the Latin for to chat, then the Latin for to say yes.
\end{quote}

\subsection*{You'll want to know: mīrābilis, mīrābile}
\textbf{mīrābilis, mīrābile} — \textquotedblleft{}wonderful, astonishing\textquotedblright{}.

\begin{quote}
Fābula mīrābilis erat, sed vēra.
\end{quote}

\begin{quote}
\emph{The story was astonishing, but true.}
\end{quote}

Two endings, from a verb meaning 'to wonder at'. English \emph{admirable} is a cousin.

\begin{cousinweb}
\textbf{mīrābilis} is built on \emph{mīror}, \textquotedblleft{}to wonder at\textquotedblright{}, from \emph{mīrus}, \textquotedblleft{}wonderful\textquotedblright{} (Lewis \& Short), with the ending \emph{-bilis}. English \emph{marvel} and \emph{miracle} come from the same family.
\end{cousinweb}

\begin{grammarlens}[title={What you've built}]
One more word for telling, explaining and describing at B1.
\end{grammarlens}

\subsection*{Your turn}
\begin{itemize}
\raggedright
  \item \emph{Say it:} \emph{mīrābilis, mīrābile}
  \item \emph{Say it:} \emph{mīrābilis, mīrābile}, once more
  \item \emph{Say it:} a sentence of your own with \emph{mīrābilis, mīrābile}
  \item \emph{Recall:} say the Latin for to chat, then the Latin for to say yes, then say \emph{mīrābilis, mīrābile} again
\end{itemize}

\begin{tcolorbox}[breakable,colback=teal!4,colframe=teal!35!black,title={Before you move on}]
What does mīrābilis, mīrābile mean? (\textquotedblleft{}Wonderful, astonishing\textquotedblright{}.) Say it once more.
\end{tcolorbox}
\section[incrēdibilis, incrēdibile]{incrēdibilis, incrēdibile — unbelievable, incredible}

\label{lesson:LA-C349-incredibilis}



\begin{quote}
Before the new one: say the Latin for to say yes, then the Latin for wonderful.
\end{quote}

\subsection*{You'll want to know: incrēdibilis, incrēdibile}
\textbf{incrēdibilis, incrēdibile} — \textquotedblleft{}unbelievable, incredible\textquotedblright{}.

\begin{quote}
Nauta rēs incrēdibilēs nārrābat.
\end{quote}

\begin{quote}
\emph{The sailor was telling unbelievable things.}
\end{quote}

Two endings: \emph{in-} 'not' + \emph{crēdibilis} 'believable', from \emph{crēdō} (taught). English \emph{incredible}.

\begin{cousinweb}
\textbf{incrēdibilis} joins \emph{in-}, \textquotedblleft{}not\textquotedblright{}, and \emph{crēdibilis}, \textquotedblleft{}worthy of belief\textquotedblright{}, from \emph{crēdō}, \textquotedblleft{}to believe\textquotedblright{} (EtymDB). English \emph{incredible} comes from it.
\end{cousinweb}

\begin{grammarlens}[title={What you've built}]
One more word for telling, explaining and describing at B1.
\end{grammarlens}

\subsection*{Your turn}
\begin{itemize}
\raggedright
  \item \emph{Say it:} \emph{incrēdibilis, incrēdibile}
  \item \emph{Say it:} \emph{incrēdibilis, incrēdibile}, once more
  \item \emph{Say it:} a sentence of your own with \emph{incrēdibilis, incrēdibile}
  \item \emph{Recall:} say the Latin for to say yes, then the Latin for wonderful, then say \emph{incrēdibilis, incrēdibile} again
\end{itemize}

\begin{tcolorbox}[breakable,colback=teal!4,colframe=teal!35!black,title={Before you move on}]
What does incrēdibilis, incrēdibile mean? (\textquotedblleft{}Unbelievable, incredible\textquotedblright{}.) Say it once more.
\end{tcolorbox}
\section[commentum, commentī]{commentum, commentī — an invention, a made-up story}

\label{lesson:LA-C349-commentum}



\begin{quote}
Before the new one: say the Latin for wonderful, then the Latin for unbelievable.
\end{quote}

\subsection*{You'll want to know: commentum, commentī}
\textbf{commentum, commentī} — \textquotedblleft{}an invention, a made-up story\textquotedblright{}.

\begin{quote}
Tōta fābula commentum erat.
\end{quote}

\begin{quote}
\emph{The whole story was an invention.}
\end{quote}

Neuter, the participle of \emph{comminīscor}, 'to think up'. Also a commentary, which gives English \emph{comment}.

\begin{cousinweb}
\textbf{commentum} is the participle of \emph{comminīscor}, \textquotedblleft{}to devise, to invent\textquotedblright{} (EtymDB, Lewis \& Short), from the root of \emph{mēns}, \textquotedblleft{}mind\textquotedblright{}. English \emph{comment} comes from it.
\end{cousinweb}

\begin{grammarlens}[title={What you've built}]
One more word for telling, explaining and describing at B1.
\end{grammarlens}

\subsection*{Your turn}
\begin{itemize}
\raggedright
  \item \emph{Say it:} \emph{commentum, commentī}
  \item \emph{Say it:} \emph{commentum, commentī}, once more
  \item \emph{Say it:} a sentence of your own with \emph{commentum, commentī}
  \item \emph{Recall:} say the Latin for wonderful, then the Latin for unbelievable, then say \emph{commentum, commentī} again
\end{itemize}

\begin{tcolorbox}[breakable,colback=teal!4,colframe=teal!35!black,title={Before you move on}]
What does commentum, commentī mean? (\textquotedblleft{}An invention, a made-up story\textquotedblright{}.) Say it once more.
\end{tcolorbox}
